\section{Conclusion}
\label{sec:conclusion}

We evaluated recursive bisection (RB) vs.\ direct n-way partitioning
across 450 chamber-year combinations (congressional, state senate, and
state house, across 50 states and 3 census years) at two geographic
resolutions, focusing primary statistical inference on the 50 independent
congressional chambers from the 2020 Census.

The key findings are:

\begin{enumerate}
  \item \textbf{RB wins on compactness across all chamber types.}
    The advantage is $+0.003$--$+0.004$ mean Polsby--Popper, statistically
    significant ($p < 0.001$) and consistent from $k = 2$ through $k = 400$.
  \item \textbf{N-way is faster only at large $k$.}
    For $k > 80$ (large state house chambers), n-way is 18--34\% faster,
    but the absolute difference is under 200\,ms --- negligible in practice.
  \item \textbf{Prime $k$ narrows the RB advantage.}
    The asymmetric bisection tree for prime $k$ loses $\sim$40\% of the
    compactness advantage.  A simple FM post-processing step recovers $\sim$60\%
    of the loss.
  \item \textbf{Block-group resolution does not change the conclusion.}
    The RB advantage scales consistently to $3\times$ larger graphs.
\end{enumerate}

These results generalise the equivalence finding of B.3 and B.4 and
provide the empirical basis for the DIA's choice of recursive bisection
as the canonical partitioning strategy for all chamber types.

\paragraph{Future work.}
An open question is whether the RB vs.\ n-way tradeoff changes under
the additional constraints of VRA compliance (B.14) or county-sticky
weights (B.10).  Preliminary results suggest the advantage persists, but
a systematic multi-constraint evaluation remains to be done.
